\section{Ten Ideas}
\label{sec:1.4}

This section summarizes some key ideas that underlie the design of many Monte Carlo methods.

\paragraph{Idea 1: Deterministic Transformations}
To sample from a given target distribution, deterministic transformation methods rely on a transport
map to convert draws from an easy-to-sample proposal distribution into draws from the target
(Chapter~2). Under mild assumptions on the target and proposal, many such transport maps exist.
Deterministic transformation methods effectively replace the problem of sampling from the target
with the problem of parameterizing and learning a transport map between the proposal and target
distributions.

\paragraph{Idea 2: Probabilistic Accept/Reject Mechanisms}
In their simplest form, probabilistic accept/reject methods rely on a probabilistic mechanism to
convert draws from a proposal distribution into draws from a given target (Chapter~2). A related
idea underlies Markov chain Monte Carlo methods (Chapter~4), where a probabilistic accept/reject
mechanism transforms samples from a proposal Markov kernel into samples from a kernel that satisfies
detailed balance with respect to the target. The unifying idea is to leverage a probabilistic rule
to turn easy-to-generate samples into approximate samples from the target distribution.

\paragraph{Idea 3: Weighting Samples}
Rather than deterministically transforming or probabilistically accepting/rejecting samples, some
Monte Carlo methods assign weights to draws from a proposal distribution to approximate expected
values with respect to the target (Chapters~3 and~9). Although this approach avoids constructing
transport maps or rejecting samples, it comes with an important limitation: the variance of the
weights can become large when target and proposal are far apart. In such cases, a small number of
samples may receive disproportionately large weights, leading to low effective sample size. This
phenomenon, known as weight degeneracy, is especially pronounced in high-dimensional settings.

\paragraph{Idea 4: Markov Chains and Local Moves}
Some of the most widely used Monte Carlo methods rely on random samples that are not independent,
but that instead form a Markov chain: the distribution of each new random draw depends on the value
of the previous one (Chapters~4, 5, 6, and~8). Positive correlation between samples increases the
variance of Monte Carlo estimates (Appendix~A), but the flexibility of the Metropolis Hastings
framework for Markov chain Monte Carlo makes up for this caveat (Chapter~4). In particular, proposal
Markov kernels that exploit information about the target distribution and its gradient can accelerate
convergence in high-dimensional settings.

\paragraph{Idea 5: Discretization of (Stochastic) Differential Equations}
A key idea in designing proposal kernels for Markov chain Monte Carlo algorithms is to leverage
Langevin stochastic differential equations and Hamiltonian differential equations that preserve the
target (Chapters~6 and~8). Time discretizations of these continuous-time dynamics yield natural
proposal kernels, which can be, if desired, combined with an accept/reject correction to ensure
convergence to the target.

\paragraph{Idea 6: Optimization}
Optimization and sampling share deep conceptual similarities, and their interplay can be leveraged
for algorithmic design. First, methods for sampling log-concave target distributions are closely
related to convex optimization algorithms: while gradient descent optimization methods can be viewed
as time-discretizations of gradient systems in parameter space, Langevin Monte Carlo sampling
methods can be viewed as discretizations of gradient flows in the space of probability densities
(Chapter~6). Second, annealing algorithms for global optimization operate by sampling from a
sequence of target distributions that become increasingly concentrated around their shared mode,
thus leveraging sampling for optimization (Chapter~7). Finally, variational inference methods seek
the closest tractable distribution to a given target, thus solving an optimization problem over
densities to approximate integrals with respect to the target (Chapter~10).

\paragraph{Idea 7: Annealing and Tempering}
To facilitate sampling from multi-modal distributions, a powerful idea is to introduce auxiliary
tempered distributions. A popular strategy in this direction is parallel tempering, where we run
several Markov chains at different temperatures and occasionally swap their states (Chapter~7).
Tempering is useful for instance in Bayesian inference, where instead of targeting the posterior
with the full data directly, a natural idea is to consider a sequence of intermediate tempered
posteriors based on smaller batches of data.

\paragraph{Idea 8: Auxiliary Variables}
Many Monte Carlo methods accelerate convergence by introducing auxiliary variables. In this
textbook, you will study several algorithms that build on this idea, including simulated tempering
(Chapter~7), Hamiltonian Monte Carlo (Chapter~8), and the slice sampler (Chapter~5). Auxiliary
variables also play a central role in sequential Monte Carlo methods (Chapter~9).

\paragraph{Idea 9: Coordinate-Wise Updates}
The Gibbs sampler and the coordinate ascent variational inference algorithm are coordinate-wise
sampling and optimization algorithms, where joint distributions are iteratively sampled or
approximated by updating one coordinate at a time (Chapters~5 and~10). This coordinate-wise
structure makes these algorithms well suited for high-dimensional problems. Gibbs samplers are
extremely effective for sampling from intractable joint distributions with easy-to-sample
conditionals. Coordinate ascent variational inference is particularly appealing in exponential
family models, where closed form formulas can be used to update the parameters.

\paragraph{Idea 10: Minimization-Maximization}
This textbook also covers the expectation maximization algorithm (Chapter~10), a prototypical
minimization-maximization method. It plays a role in maximum likelihood estimation analogous to
that of the Gibbs sampler in posterior sampling. In particular, the expectation maximization
algorithm can be used to provide effective initializations for Gibbs samplers.
